% 14. Organizacja pracy
\section{Organizacja pracy}\label{sec:organizacja}

\begin{instrukcja}[wymagane]
Kto za co odpowiadał i jak pracowaliście. Każda osoba musi umieć wyjaśnić dowolny fragment kodu -
także ten, którego nie pisała. Przy narzędziach AI napiszcie konkretnie, do czego ich używaliście
i czego nie robiły za Was.
\end{instrukcja}

\subsection{Podział pracy}
\begin{table}[H]\centering
\caption{Podział pracy w zespole}\label{tab:podzial}
\begin{tabularx}{\linewidth}{@{}p{4.4cm}Lp{2cm}@{}}
\naglowektabeli \nt{Osoba} & \nt{Za co odpowiadała} & \nt{Udział} \\
Mateusz Machowski & ramy graficzne, nawigacja, ekrany E1-E3, wąski ekran; rozdziały 4, 5, 8 & \wpisz{50\%} \\
Bartosz Kurbiel & repozytorium, formularz E4, walidacja, dane i operacje; rozdziały 7, 9 & \wpisz{50\%} \\
\bottomrule
\end{tabularx}
\end{table}

\subsection{Sposób pracy}
Zadania podzieliliśmy według terminów z harmonogramu w pliku \texttt{docs/PODZIAL\_PRACY.md}: osoba A
odpowiada za ramy, start, listę i szczegóły, osoba B za formularz, dane i operacje. Kod jest w jednym
repozytorium na GitHubie; każdą zmianę ogląda druga osoba.

\subsection{Wykorzystanie narzędzi AI}
\begin{tabularx}{\linewidth}{@{}p{3.6cm}p{3.2cm}L@{}}
\naglowektabeli \nt{Narzędzie} & \nt{Kto używał} & \nt{Do czego} \\
Claude Code (Anthropic) & M. Machowski & szkic komponentów listy, karty i szczegółów, efektów wizualnych (licznik dni, podświetlenie kart), zrzuty ekranu i pierwsza wersja tekstu dokumentacji \\
\wpisz{...} & & \\
\bottomrule
\end{tabularx}

\medskip
Decyzje projektowe (model nawigacji, wybór wzorców, zakres) podjęliśmy sami na podstawie wykładów
i makiet z K1. Wygenerowany kod przeglądaliśmy i uruchamialiśmy na szerokim i wąskim ekranie; podczas
sprawdzania zrzutów wyszły dwa błędy na telefonie (pasek menu przykryty przez karty, za szerokie pole
wyboru), które poprawiliśmy. Każda osoba potrafi wyjaśnić kod obu części.
